\documentclass[10pt,a4paper]{amsart}
\usepackage{amsmath,amssymb,amsthm,graphicx}
% Retained preamble for the Stacks Project chapter {CH} (shared project preamble).
% The chapter file pinned-original-source.tex only carries \input{{preamble}} on its first line,
% so no preamble text lives inside the chapter itself.  The Stacks Project uses one single shared
% preamble for every chapter: the snapshot stages it three times (sources/stacks-chow/preamble.tex,
% sources/stacks-etale-cohomology/preamble.tex and focused-seeds/AGP/sources/stacks-snapshot-20261011/
% preamble.tex) and those copies are identical in size.  The block below is copied mechanically,
% line for line, from sources/stacks-chow/preamble.tex lines 14-18 and 154-217: the Xy-pic setup and
% the native theorem declarations together with the mathematical macros.  Nothing was added, removed
% or redefined; the wrapper class and the amsmath/amssymb/amsthm/graphicx base packages are supplied
% by the assembler.
% For commutative diagrams we use Xy-pic
% Chapter: stacks-cohomology.tex -- Cohomology of Sheaves.
\usepackage[all]{xy}

% We use 2cell for 2-commutative diagrams.
\xyoption{2cell}

% Theorem environments.
%
\theoremstyle{plain}
\newtheorem{theorem}[subsection]{Theorem}
\newtheorem{proposition}[subsection]{Proposition}
\newtheorem{lemma}[subsection]{Lemma}

\theoremstyle{definition}
\newtheorem{definition}[subsection]{Definition}
\newtheorem{example}[subsection]{Example}
\newtheorem{exercise}[subsection]{Exercise}
\newtheorem{situation}[subsection]{Situation}

\theoremstyle{remark}
\newtheorem{remark}[subsection]{Remark}
\newtheorem{remarks}[subsection]{Remarks}

\numberwithin{equation}{subsection}

% Macros
%
\def\lim{\mathop{\mathrm{lim}}\nolimits}
\def\colim{\mathop{\mathrm{colim}}\nolimits}
\def\Spec{\mathop{\mathrm{Spec}}}
\def\Hom{\mathop{\mathrm{Hom}}\nolimits}
\def\Ext{\mathop{\mathrm{Ext}}\nolimits}
\def\SheafHom{\mathop{\mathcal{H}\!\mathit{om}}\nolimits}
\def\SheafExt{\mathop{\mathcal{E}\!\mathit{xt}}\nolimits}
\def\Sch{\mathit{Sch}}
\def\Mor{\mathop{\mathrm{Mor}}\nolimits}
\def\Ob{\mathop{\mathrm{Ob}}\nolimits}
\def\Sh{\mathop{\mathit{Sh}}\nolimits}
\def\NL{\mathop{N\!L}\nolimits}
\def\CH{\mathop{\mathrm{CH}}\nolimits}
\def\proetale{{pro\text{-}\acute{e}tale}}
\def\etale{{\acute{e}tale}}
\def\QCoh{\mathit{QCoh}}
\def\Ker{\mathop{\mathrm{Ker}}}
\def\Im{\mathop{\mathrm{Im}}}
\def\Coker{\mathop{\mathrm{Coker}}}
\def\Coim{\mathop{\mathrm{Coim}}}

% Boxtimes
%
\DeclareMathSymbol{\boxtimes}{\mathbin}{AMSa}{"02}

%
% Macros for moduli stacks/spaces
%
\def\QCohstack{\mathcal{QC}\!\mathit{oh}}
\def\Cohstack{\mathcal{C}\!\mathit{oh}}
\def\Spacesstack{\mathcal{S}\!\mathit{paces}}
\def\Quotfunctor{\mathrm{Quot}}
\def\Hilbfunctor{\mathrm{Hilb}}
\def\Curvesstack{\mathcal{C}\!\mathit{urves}}
\def\Polarizedstack{\mathcal{P}\!\mathit{olarized}}
\def\Complexesstack{\mathcal{C}\!\mathit{omplexes}}
% \Pic is the operator that assigns to X its picard group, usage \Pic(X)
% \Picardstack_{X/B} denotes the Picard stack of X over B
% \Picardfunctor_{X/B} denotes the Picard functor of X over B
\def\Pic{\mathop{\mathrm{Pic}}\nolimits}
\def\Picardstack{\mathcal{P}\!\mathit{ic}}
\def\Picardfunctor{\mathrm{Pic}}
\def\Deformationcategory{\mathcal{D}\!\mathit{ef}}
\usepackage{cleveref}
\crefname{originallemma}{original lemma}{original lemmas}
\Crefname{originallemma}{Original lemma}{Original lemmas}
\makeatletter
\newlabel{derived-lemma-leray-acyclicity}{{015E}{source}{}{sourcepointer-1}{}}
\newlabel{derived-lemma-leray-acyclicity@cref}{{[originallemma][0][]015E}{[1][0][]source}{}{}{}}
\makeatother
\crefname{originalequation}{original equation}{original equations}
\Crefname{originalequation}{Original equation}{Original equations}
\makeatletter
\newlabel{equation-needs-signs}{{07MB}{source}{}{sourcepointer-2}{}}
\newlabel{equation-needs-signs@cref}{{[originalequation][0][]07MB}{[1][0][]source}{}{}{}}
\makeatother
\crefname{originallemma}{original lemma}{original lemmas}
\Crefname{originallemma}{Original lemma}{Original lemmas}
\makeatletter
\newlabel{lemma-cech-complex-complex}{{08BN}{source}{}{sourcepointer-3}{}}
\newlabel{lemma-cech-complex-complex@cref}{{[originallemma][0][]08BN}{[1][0][]source}{}{}{}}
\makeatother
\crefname{originallemma}{original lemma}{original lemmas}
\Crefname{originallemma}{Original lemma}{Original lemmas}
\makeatletter
\newlabel{lemma-cech-complex-complex-computes}{{0FLH}{source}{}{sourcepointer-4}{}}
\newlabel{lemma-cech-complex-complex-computes@cref}{{[originallemma][0][]0FLH}{[1][0][]source}{}{}{}}
\makeatother
\crefname{originallemma}{original lemma}{original lemmas}
\Crefname{originallemma}{Original lemma}{Original lemmas}
\makeatletter
\newlabel{lemma-cup-compatible-with-naive}{{0FP3}{source}{}{sourcepointer-5}{}}
\newlabel{lemma-cup-compatible-with-naive@cref}{{[originallemma][0][]0FP3}{[1][0][]source}{}{}{}}
\makeatother
\crefname{originallemma}{original lemma}{original lemmas}
\Crefname{originallemma}{Original lemma}{Original lemmas}
\makeatletter
\newlabel{lemma-flasque-acyclic}{{09SY}{source}{}{sourcepointer-6}{}}
\newlabel{lemma-flasque-acyclic@cref}{{[originallemma][0][]09SY}{[1][0][]source}{}{}{}}
\makeatother
\crefname{originallemma}{original lemma}{original lemmas}
\Crefname{originallemma}{Original lemma}{Original lemmas}
\makeatletter
\newlabel{lemma-godement-resolution-bounded-below}{{0FKT}{source}{}{sourcepointer-7}{}}
\newlabel{lemma-godement-resolution-bounded-below@cref}{{[originallemma][0][]0FKT}{[1][0][]source}{}{}{}}
\makeatother
\crefname{originallemma}{original lemma}{original lemmas}
\Crefname{originallemma}{Original lemma}{Original lemmas}
\makeatletter
\newlabel{more-algebra-lemma-derived-tor-homotopy}{{064I}{source}{}{sourcepointer-8}{}}
\newlabel{more-algebra-lemma-derived-tor-homotopy@cref}{{[originallemma][0][]064I}{[1][0][]source}{}{}{}}
\makeatother
\title{Commutativity of the diagram comparing the naive cup product with the derived cup product}
\author{The Stacks Project Authors}
\date{stacks-project:cohomology.tex:local-snapshot-20261011; GNU FDL 1.2 or later}
\begin{document}
\maketitle
\noindent\textit{Editorial note.} All delivered mathematical text is the author's own, copied from the pinned Stacks Project chapter file cohomology.tex; no statement or proof step was rewritten, reordered or removed.\par\medskip
\noindent\textit{Editorial note.} The statement of the lemma contains the label equation-needs-signs inside a xymatrix, which is a pointer to an equation outside the counted span and is carried as a navigation stub; the equation itself is not load bearing for the counted argument.\par\medskip
\noindent\textit{Editorial note.} The counted spans carry 12 navigation stubs and no bibliography citation.\par\medskip
\noindent\textit{Editorial note.} The preamble is the chapter's retained minimal preamble (sources/stacks-cohomology/retained-preamble.tex lines 1-81), a verbatim excerpt of the native preamble; the Xy-pic setup it contains is what renders the diagrams of the statement and the proof.\par\medskip
\noindent\textit{Editorial note.} The label lemma-diagrams-commute is defined only in cohomology.tex, so no cross-chapter restatement had to be weighed for this card.\par\medskip
\noindent\textit{Allowed standard prerequisites (editorial).}\begin{itemize}
\item Godement resolution and its quasi-isomorphism property. Same chapter, Lemma lemma-godement-resolution-bounded-below, whose quasi-isomorphisms a and b are chosen at proof lines 7581-7584. Supplies the replacement complexes K\_1 and M\_1 used in the first reduction.
\item A stalkwise homotopy equivalence of complexes of modules induces a quasi-isomorphism on total tensor products. More on Algebra, Lemma more-algebra-lemma-derived-tor-homotopy, invoked at proof lines 7592-7593. Supplies that the targets of the two diagrams agree in D(A) after the replacement.
\item Flasque acyclicity and the Leray acyclicity lemma for the global sections functor. Same chapter, Lemma lemma-flasque-acyclic and Derived Categories, Lemma derived-lemma-leray-acyclicity, invoked at proof lines 7648-7651. Supplies that the horizontal arrows of the comparison diagram are isomorphisms in D(A).
\item Compatibility of the naive cup product with the derived cup product. Same chapter, Lemma lemma-cup-compatible-with-naive, invoked at proof lines 7652-7655. Supplies the commutativity statement that is transferred to the diagram of the lemma.
\end{itemize}\medskip
\ifdefined\hypertarget\hypertarget{sourcepointer-1}{}\fi
\noindent\textit{Source reference.} Original lemma 015E: sources/stacks-chow/tags.txt, source lines 1505--1505. Pointer outside the counted span, invoked once for context; the load-bearing argument of the counted statement is the author's proof reproduced in full in the proof span.\par\medskip
\ifdefined\hypertarget\hypertarget{sourcepointer-2}{}\fi
\noindent\textit{Source reference.} Original equation 07MB: sources/stacks-cohomology/pinned-original-source.tex, source lines 5273--5273. Pointer outside the counted span, invoked once for context; the load-bearing argument of the counted statement is the author's proof reproduced in full in the proof span.\par\medskip
\ifdefined\hypertarget\hypertarget{sourcepointer-3}{}\fi
\noindent\textit{Source reference.} Original lemma 08BN: sources/stacks-cohomology/pinned-original-source.tex, source lines 4939--4939. Pointer outside the counted span, invoked once for context; the load-bearing argument of the counted statement is the author's proof reproduced in full in the proof span.\par\medskip
\ifdefined\hypertarget\hypertarget{sourcepointer-4}{}\fi
\noindent\textit{Source reference.} Original lemma 0FLH: sources/stacks-cohomology/pinned-original-source.tex, source lines 5022--5022. Pointer outside the counted span, invoked once for context; the load-bearing argument of the counted statement is the author's proof reproduced in full in the proof span.\par\medskip
\ifdefined\hypertarget\hypertarget{sourcepointer-5}{}\fi
\noindent\textit{Source reference.} Original lemma 0FP3: sources/stacks-cohomology/pinned-original-source.tex, source lines 7312--7312. Pointer outside the counted span, invoked once for context; the load-bearing argument of the counted statement is the author's proof reproduced in full in the proof span.\par\medskip
\ifdefined\hypertarget\hypertarget{sourcepointer-6}{}\fi
\noindent\textit{Source reference.} Original lemma 09SY: sources/stacks-cohomology/pinned-original-source.tex, source lines 1913--1913. Pointer outside the counted span, invoked once for context; the load-bearing argument of the counted statement is the author's proof reproduced in full in the proof span.\par\medskip
\ifdefined\hypertarget\hypertarget{sourcepointer-7}{}\fi
\noindent\textit{Source reference.} Original lemma 0FKT: sources/stacks-cohomology/pinned-original-source.tex, source lines 7081--7081. Pointer outside the counted span, invoked once for context; the load-bearing argument of the counted statement is the author's proof reproduced in full in the proof span.\par\medskip
\ifdefined\hypertarget\hypertarget{sourcepointer-8}{}\fi
\noindent\textit{Source reference.} Original lemma 064I: sources/stacks-chow/tags.txt, source lines 7867--7867. Pointer outside the counted span, invoked once for context; the load-bearing argument of the counted statement is the author's proof reproduced in full in the proof span.\par\medskip
\begin{lemma}
\label{lemma-diagrams-commute}
Let $(X, \mathcal{O}_X)$ be a ringed space. Let
$\mathcal{K}^\bullet$ and $\mathcal{M}^\bullet$
be bounded below complexes of $\mathcal{O}_X$-modules.
Let $\mathcal{U} : X = \bigcup_{i \in I} U_i$ be an open covering
Then
$$
\xymatrix{
\text{Tot}(\check{\mathcal{C}}^\bullet(\mathcal{U}, \mathcal{K}^\bullet))
\otimes_A^\mathbf{L}
\text{Tot}(\check{\mathcal{C}}^\bullet(\mathcal{U}, \mathcal{M}^\bullet))
\ar[d] \ar[r] &
R\Gamma(X, \mathcal{K}^\bullet)
\otimes_A^\mathbf{L}
R\Gamma(X, \mathcal{M}^\bullet) \ar[d]^\mu \\
\text{Tot}(
\text{Tot}(\check{\mathcal{C}}^\bullet(\mathcal{U}, \mathcal{K}^\bullet))
\otimes_A
\text{Tot}(\check{\mathcal{C}}^\bullet(\mathcal{U}, \mathcal{M}^\bullet)))
\ar[d]^{(\ref{equation-needs-signs})} &
R\Gamma(X,
\mathcal{K}^\bullet \otimes_{\mathcal{O}_X}^\mathbf{L} \mathcal{M}^\bullet)
\ar[d] \\
\text{Tot}(
\check{\mathcal{C}}^\bullet({\mathcal U},
\text{Tot}(\mathcal{K}^\bullet \otimes_{\mathcal{O}_X} \mathcal{M}^\bullet)
)) \ar[r] &
R\Gamma(X,
\text{Tot}(\mathcal{K}^\bullet \otimes_{\mathcal{O}_X} \mathcal{M}^\bullet))
}
$$
where the horizontal arrows are the ones in
Lemma \ref{lemma-cech-complex-complex}
commutes in $D(A)$.
\end{lemma}

\begin{proof}
Choose quasi-isomorphisms of complexes
$a : \mathcal{K}^\bullet \to \mathcal{K}_1^\bullet$ and
$b : \mathcal{M}^\bullet \to \mathcal{M}_1^\bullet$
as in Lemma \ref{lemma-godement-resolution-bounded-below}.
Since the maps $a$ and $b$ on stalks are homotopy equivalences
we see that the induced map
$$
\text{Tot}(\mathcal{K}^\bullet \otimes_{\mathcal{O}_X} \mathcal{M}^\bullet)
\to
\text{Tot}(\mathcal{K}_1^\bullet \otimes_{\mathcal{O}_X} \mathcal{M}_1^\bullet)
$$
is a homotopy equivalence on stalks too (More on Algebra, Lemma
\ref{more-algebra-lemma-derived-tor-homotopy}) and hence a quasi-isomorphism.
Thus the targets
$$
R\Gamma(X,
\text{Tot}(\mathcal{K}^\bullet
\otimes_{\mathcal{O}_X} \mathcal{M}^\bullet)) =
R\Gamma(X,
\text{Tot}(\mathcal{K}_1^\bullet
\otimes_{\mathcal{O}_X} \mathcal{M}_1^\bullet))
$$
of the two diagrams are the same in $D(A)$. It follows that it suffices
to prove the diagram commutes for $\mathcal{K}$ and $\mathcal{M}$
replaced by $\mathcal{K}_1$ and $\mathcal{M}_1$. This reduces us to
the case discussed in the next paragraph.

\medskip\noindent
Assume $\mathcal{K}^\bullet$ and $\mathcal{M}^\bullet$ are bounded
below complexes of flasque $\mathcal{O}_X$-modules and
consider the diagram relating the cup product with the cup product
(\ref{equation-needs-signs}) on {\v C}ech complexes.
Then we can consider the commutative diagram
$$
\xymatrix{
\Gamma(X, \mathcal{K}^\bullet)
\otimes_A^\mathbf{L}
\Gamma(X, \mathcal{M}^\bullet) \ar[d] \ar[r] &
\text{Tot}(\check{\mathcal{C}}^\bullet(\mathcal{U}, \mathcal{K}^\bullet))
\otimes_A^\mathbf{L}
\text{Tot}(\check{\mathcal{C}}^\bullet(\mathcal{U}, \mathcal{M}^\bullet))
\ar[d] \\
\text{Tot}(\Gamma(X, \mathcal{K}^\bullet)
\otimes_A
\Gamma(X, \mathcal{M}^\bullet)) \ar[d] \ar[r] &
\text{Tot}(
\text{Tot}(\check{\mathcal{C}}^\bullet(\mathcal{U}, \mathcal{K}^\bullet))
\otimes_A
\text{Tot}(\check{\mathcal{C}}^\bullet(\mathcal{U}, \mathcal{M}^\bullet)))
\ar[d]^{(\ref{equation-needs-signs})} \\
\Gamma(X, \text{Tot}(\mathcal{K}^\bullet
\otimes_{\mathcal{O}_X}
\mathcal{M}^\bullet)) \ar[r] &
\text{Tot}(
\check{\mathcal{C}}^\bullet({\mathcal U},
\text{Tot}(\mathcal{K}^\bullet \otimes_{\mathcal{O}_X} \mathcal{M}^\bullet)
))
}
$$
In this diagram the horizontal arrows are isomorphisms in $D(A)$ because
for a bounded below complex of flasque modules such as $\mathcal{K}^\bullet$
we have
$$
\Gamma(X, \mathcal{K}^\bullet) =
\text{Tot}(\check{\mathcal{C}}^\bullet(\mathcal{U}, \mathcal{K}^\bullet)) =
R\Gamma(X, \mathcal{K}^\bullet)
$$
in $D(A)$. This follows from
Lemma \ref{lemma-flasque-acyclic},
Derived Categories, Lemma \ref{derived-lemma-leray-acyclicity}, and
Lemma \ref{lemma-cech-complex-complex-computes}.
Hence the commutativity of the diagram of the lemma involving
(\ref{equation-needs-signs}) follows from the already proven
commutativity of Lemma \ref{lemma-cup-compatible-with-naive}
where $f$ is the morphism to a point (see discussion
following Lemma \ref{lemma-cup-compatible-with-naive}).
\end{proof}

\end{document}
